\section{应用图谱：数字工作流如何成为 Agent 的主战场}

\subsection{任务形态：单应用操作到跨应用编排}
讲者强调，实际用户任务通常不是一个页面一次点击，而是跨工具流程。例如：浏览器检索信息、表格整理数据、邮件发送结果、文档归档记录。多模态 Agent 需要处理持续状态与任务上下文，而不是短期指令响应。

\begin{importantbox}{真实世界任务的四个特征}
\begin{enumerate}
    \item \textbf{长链路}：常常超过 20--50 步；
    \item \textbf{跨界面}：浏览器、桌面、终端、Office 工具混合出现；
    \item \textbf{非确定性}：页面加载、弹窗、网络波动导致状态变化；
    \item \textbf{目标开放}：很多任务存在多条可行路径，不是单一标准答案。
\end{enumerate}
\end{importantbox}

\subsection{Agent 能力分层：感知、规划、执行、验证}
课程把 Agent 拆成能力栈，而不是一个黑盒模型。感知层负责理解界面与语义对象；规划层决定下一步子目标；执行层输出可执行动作；验证层判定任务是否完成并修正轨迹。这种分层能更好定位系统瓶颈。

\begin{knowledgebox}{为什么多模态 Agent 需要 ``过程可见性''}
当任务失败时，团队必须知道问题在感知、规划还是执行。如果所有失败都被归因到 ``模型不够强''，就无法制定有效优化策略。可观测性是 Agent 迭代效率的核心。
\end{knowledgebox}

\subsection{数字场景优先，物理场景随后}
从落地节奏看，讲者明显把数字环境（web、desktop、mobile）作为优先级更高的方向。原因是反馈闭环更快、评测脚本更容易搭建、成本更可控。物理世界 Agent 虽然重要，但数据与评测基础设施更重，迭代速度更慢。

\begin{warningbox}{不要过早把机器人路径与 GUI Agent 混为一谈}
两者都属于 Agent，但工程约束完全不同。GUI Agent 的瓶颈在界面理解和任务编排；物理 Agent 还叠加感知误差、安全约束和执行器稳定性。混合讨论会掩盖问题本质。
\end{warningbox}

\subsection{本章小结}
课程把应用问题定义为 ``跨应用、多步骤、可恢复'' 的长期任务。由此引申出的工程要求是：先把数字任务闭环做深，再向更复杂场景扩展。

